\documentclass{article}
\usepackage{amsmath}
\usepackage{amsfonts}
\usepackage{amssymb}
\usepackage{graphicx}
\usepackage[usenames,dvipsnames]{xcolor}
\usepackage{hyperref}
\usepackage{booktabs}
\usepackage{array}

\pagestyle{plain}
\textheight235truemm\textwidth165truemm
\voffset-2.2cm\hoffset-1.2cm
\vfuzz=10000truept\hfuzz=10000truept

\newcommand{\RED}[1]{{\color{red}{#1}}}
\newcommand{\R}{{\mathbb{R}}}
\newtheorem{theorem}{Theorem}[section]
\newtheorem{lemma}[theorem]{Lemma}
\newtheorem{definition}[theorem]{Definition}
\newtheorem{remark}[theorem]{Remark}

\usepackage[
backend=biber,
style=alphabetic,
sorting=ynt
]{biblatex}
\addbibresource{Chinese_med.bib}

\title{Dynamics of Five Organ Diagram in Chinese Medicine}
\author{Carter Feng and Yunjiao Wang}
\date{September 2025}

\begin{document}

\maketitle

%===========================================================
\section{Relations in Wu-Xing and Mathematical Modeling}
%===========================================================

\subsection{The Five Organs and Index Assignment}

We index the five organs as follows, following the standard Wu-Xing (Five Element)
cycle:

\begin{center}
\begin{tabular}{cll}
\toprule
Index & Element & Organ \\
\midrule
$x_1$ & Wood\; (木) & Liver\; (肝) \\
$x_2$ & Fire\; (火)  & Heart\; (心) \\
$x_3$ & Earth\; (土) & Spleen\; (脾) \\
$x_4$ & Metal\; (金) & Lung\; (肺) \\
$x_5$ & Water\; (水) & Kidney\; (肾) \\
\bottomrule
\end{tabular}
\end{center}

\noindent
Let $x_i = x_i(t) \geq 0$ denote the functional activity level (qi) of organ $i$
at time $t$.

%-----------------------------------------------------------
\subsection{Fundamental Relationships}
%-----------------------------------------------------------

TCM identifies three primary inter-organ relationships, each of which contributes
a term to the dynamical model.

\medskip
\noindent\textbf{Sheng (生, Generating) cycle.}
Each organ nourishes the next in the cycle
\[
x_1 \to x_2 \to x_3 \to x_4 \to x_5 \to x_1.
\]
We denote the \emph{Sheng mother} of organ $i$ by $\sigma(i)$:
\[
\sigma(1)=5,\quad \sigma(2)=1,\quad \sigma(3)=2,\quad \sigma(4)=3,\quad \sigma(5)=4.
\]

\medskip
\noindent\textbf{Ke (克, Controlling) cycle.}
Each organ restrains another following the pentagram pattern
\[
x_1 \to x_3 \to x_5 \to x_2 \to x_4 \to x_1.
\]
We denote the \emph{Ke target} of organ $i$ (the organ that $i$ controls) by
$\kappa(i)$, and the \emph{Ke controller} of organ $i$ (the organ that controls $i$)
by $\rho(i)$:

\begin{center}
\begin{tabular}{cll}
\toprule
Organ $i$ & Ke target $\kappa(i)$ & Ke controller $\rho(i)$ \\
\midrule
$x_1$ Wood  & $x_3$ Earth  & $x_4$ Metal  \\
$x_2$ Fire  & $x_4$ Metal  & $x_5$ Water  \\
$x_3$ Earth & $x_5$ Water  & $x_1$ Wood   \\
$x_4$ Metal & $x_1$ Wood   & $x_2$ Fire   \\
$x_5$ Water & $x_2$ Fire   & $x_3$ Earth  \\
\bottomrule
\end{tabular}
\end{center}

\medskip
\noindent\textbf{Wu (侮, Rebellious Counter-Acting).}
Under pathological conditions, when organ $i$ becomes severely deficient, the organ
it normally controls, $x_{\kappa(i)}$, may rebel against it. That is, $x_{\kappa(i)}$
exerts a reverse draining effect on $x_i$ when $x_i$ falls below a threshold
$\epsilon_i > 0$.

%-----------------------------------------------------------
\subsection{Adaptive Set Point}
%-----------------------------------------------------------

Rather than a fixed homeostatic target, we posit that each organ maintains an
\emph{adaptive set point} $\theta_i > 0$ that is influenced by its Sheng mother
and its Ke controller:
\begin{equation}\label{eq:setpoint}
    \theta_i \;=\; \theta_i^0 \;+\; \alpha_i\, x_{\sigma(i)} \;-\; \beta_i\, x_{\rho(i)},
\end{equation}
where $\theta_i^0 > 0$ is the baseline constitutional set point of organ $i$,
$\alpha_i \geq 0$ encodes how much the mother's nourishment raises the set point,
and $\beta_i \geq 0$ encodes how much the controller's restraint lowers it.
We assume parameters are chosen such that $\theta_i > 0$ for all physiologically
relevant states.

%-----------------------------------------------------------
\subsection{Model Derivation: Individual Terms}
%-----------------------------------------------------------

We now describe each term in the model.

\medskip
\noindent\textbf{Adaptive self-production (Hill-type).}
In TCM, a healthy organ maintains continuous baseline production of qi. When the
organ is in excess ($x_i > \theta_i$), production is downregulated; when the organ
is deficient ($x_i < \theta_i$), production is upregulated. We model this with
two Hill functions.

The first term models \emph{activity-driven production} with downregulation in excess:
\[
    \frac{a_i\, x_i\, \theta_i^{n_i}}{\theta_i^{n_i} + x_i^{n_i}},
\]
where $a_i > 0$ is the maximum production rate and $n_i \geq 1$ is the Hill
coefficient controlling the sigmoidal sharpness. When $x_i \ll \theta_i$ this
approaches $a_i x_i$ (full production); when $x_i \gg \theta_i$ it approaches $0$
(production shuts off). At the set point $x_i = \theta_i$ the rate equals
$\tfrac{1}{2}a_i x_i$, which is nonzero, so the organ continues to function at
its healthy level.

The second term models \emph{constitutive upregulation} in deficiency:
\[
    \frac{u_i\, \theta_i^{m_i}}{\theta_i^{m_i} + x_i^{m_i}},
\]
where $u_i > 0$ is the maximum compensatory boost and $m_i \geq 1$ is a Hill
coefficient. This term is large when $x_i$ is small and small when $x_i$ is large,
capturing the TCM principle of zheng qi (正气) compensation: a depleted organ
mobilizes additional resources to recover. The two exponents $n_i$ and $m_i$
allow the excess-downregulation and deficiency-upregulation responses to be
asymmetric.

\medskip
\noindent\textbf{Sheng nourishment.}
The Sheng mother $x_{\sigma(i)}$ nourishes organ $i$, with saturation as $x_i$
grows (a strong organ requires less maternal support):
\[
    \frac{c_i\, x_{\sigma(i)}}{1 + x_i},
\]
where $c_i > 0$ is the nourishment strength.

\medskip
\noindent\textbf{Ke over-acting (Cheng, 乘).}
Organ $i$ over-acts on its Ke target $x_{\kappa(i)}$, draining it. As the target
recovers, it increasingly resists this drain:
\[
    -\frac{g_i\, x_i}{1 + x_{\kappa(i)}},
\]
where $g_i > 0$ is the over-acting strength. The denominator $1 + x_{\kappa(i)}$
means the drain diminishes as the target strengthens.

\medskip
\noindent\textbf{Ke direct restraint.}
Organ $i$ is restrained by its Ke controller $x_{\rho(i)}$. Both organs must be
active for restraint to occur, motivating a mass-action term:
\[
    -h_i\, x_i\, x_{\rho(i)},
\]
where $h_i > 0$ is the restraint strength. Note that the Ke over-acting term and
the Ke restraint term are distinct: the former describes what organ $i$ does to
its target, while the latter describes what the controller does to organ $i$.

\medskip
\noindent\textbf{Wu rebellious counter-acting.}
When organ $i$ is severely deficient (below threshold $\epsilon_i$), its Ke target
$x_{\kappa(i)}$ rebels against it. The rebellion strength grows as $x_i$ weakens:
\[
    -\frac{w_i\, x_{\kappa(i)}}{1 + x_i/\epsilon_i},
\]
where $w_i \geq 0$ is the rebellion strength and $\epsilon_i > 0$ is the weakness
threshold. When $x_i \gg \epsilon_i$ (organ $i$ is healthy) this term vanishes;
when $x_i \ll \epsilon_i$ it approaches $-w_i x_{\kappa(i)}$, actively draining
organ $i$ in proportion to the rebel's activity. Setting $w_i = 0$ for all $i$
recovers a homeostasis-only model without pathological dynamics.

%-----------------------------------------------------------
\subsection{Full Model}
%-----------------------------------------------------------

Combining all terms, the full five-dimensional dynamical system is:

\begin{equation}\label{eq:fullmodel}
\boxed{
\dot{x}_i \;=\;
\underbrace{\frac{a_i\, x_i\, \theta_i^{n_i}}{\theta_i^{n_i} + x_i^{n_i}}}_{\text{Hill production}}
\;+\;
\underbrace{\frac{u_i\, \theta_i^{m_i}}{\theta_i^{m_i} + x_i^{m_i}}}_{\text{deficiency upregulation}}
\;+\;
\underbrace{\frac{c_i\, x_{\sigma(i)}}{1 + x_i}}_{\text{Sheng}}
\;-\;
\underbrace{\frac{g_i\, x_i}{1 + x_{\kappa(i)}}}_{\text{Ke over-acting}}
\;-\;
\underbrace{h_i\, x_i\, x_{\rho(i)}}_{\text{Ke restraint}}
\;-\;
\underbrace{\frac{w_i\, x_{\kappa(i)}}{1 + x_i/\epsilon_i}}_{\text{Wu rebellion}}
}
\end{equation}

\noindent
for $i = 1, 2, 3, 4, 5$, where the adaptive set point $\theta_i$ is given by
\eqref{eq:setpoint}. Written out explicitly for each organ:

\begin{equation}\label{eq:explicit}
\begin{aligned}
\dot{x}_1 &= \frac{a_1 x_1 \theta_1^{n_1}}{\theta_1^{n_1} + x_1^{n_1}}
           + \frac{u_1 \theta_1^{m_1}}{\theta_1^{m_1} + x_1^{m_1}}
           + \frac{c_1 x_5}{1 + x_1}
           - \frac{g_1 x_1}{1 + x_3}
           - h_1 x_1 x_4
           - \frac{w_1 x_3}{1 + x_1/\epsilon_1}\\[8pt]
\dot{x}_2 &= \frac{a_2 x_2 \theta_2^{n_2}}{\theta_2^{n_2} + x_2^{n_2}}
           + \frac{u_2 \theta_2^{m_2}}{\theta_2^{m_2} + x_2^{m_2}}
           + \frac{c_2 x_1}{1 + x_2}
           - \frac{g_2 x_2}{1 + x_4}
           - h_2 x_2 x_5
           - \frac{w_2 x_4}{1 + x_2/\epsilon_2}\\[8pt]
\dot{x}_3 &= \frac{a_3 x_3 \theta_3^{n_3}}{\theta_3^{n_3} + x_3^{n_3}}
           + \frac{u_3 \theta_3^{m_3}}{\theta_3^{m_3} + x_3^{m_3}}
           + \frac{c_3 x_2}{1 + x_3}
           - \frac{g_3 x_3}{1 + x_5}
           - h_3 x_3 x_1
           - \frac{w_3 x_5}{1 + x_3/\epsilon_3}\\[8pt]
\dot{x}_4 &= \frac{a_4 x_4 \theta_4^{n_4}}{\theta_4^{n_4} + x_4^{n_4}}
           + \frac{u_4 \theta_4^{m_4}}{\theta_4^{m_4} + x_4^{m_4}}
           + \frac{c_4 x_3}{1 + x_4}
           - \frac{g_4 x_4}{1 + x_1}
           - h_4 x_4 x_2
           - \frac{w_4 x_1}{1 + x_4/\epsilon_4}\\[8pt]
\dot{x}_5 &= \frac{a_5 x_5 \theta_5^{n_5}}{\theta_5^{n_5} + x_5^{n_5}}
           + \frac{u_5 \theta_5^{m_5}}{\theta_5^{m_5} + x_5^{m_5}}
           + \frac{c_5 x_4}{1 + x_5}
           - \frac{g_5 x_5}{1 + x_2}
           - h_5 x_5 x_3
           - \frac{w_5 x_2}{1 + x_5/\epsilon_5}
\end{aligned}
\end{equation}

\noindent
with $\theta_i = \theta_i^0 + \alpha_i x_{\sigma(i)} - \beta_i x_{\rho(i)}$
substituted inline.

%-----------------------------------------------------------
\subsection{Parameter Summary}
%-----------------------------------------------------------

\begin{center}
\begin{tabular}{cll}
\toprule
Parameter & Biological meaning & Count \\
\midrule
$\theta_i^0$ & Baseline constitutional set point of organ $i$ & 5 \\
$\alpha_i$   & Sensitivity of $\theta_i$ to Sheng mother nourishment & 5 \\
$\beta_i$    & Sensitivity of $\theta_i$ to Ke controller restraint & 5 \\
$a_i$        & Maximum self-production rate & 5 \\
$n_i$        & Hill exponent for production downregulation in excess & 5 \\
$u_i$        & Maximum compensatory upregulation in deficiency & 5 \\
$m_i$        & Hill exponent for upregulation in deficiency & 5 \\
$c_i$        & Sheng nourishment strength & 5 \\
$g_i$        & Ke over-acting (Cheng) strength & 5 \\
$h_i$        & Ke direct restraint strength & 5 \\
$w_i$        & Wu rebellion strength ($w_i = 0$: homeostasis only) & 5 \\
$\epsilon_i$ & Weakness threshold for Wu rebellion activation & 5 \\
\midrule
 & Total & 60 \\
\bottomrule
\end{tabular}
\end{center}

\noindent
For a simplified symmetric analysis one may set $n_i = m_i = n$ globally,
reducing the parameter count to 56. Further reduction to a symmetric system
($a_i = a$, $c_i = c$, etc.) yields a tractable starting point for equilibrium
and bifurcation analysis, with $n$ as a natural bifurcation parameter governing
the sharpness of homeostatic response.

\begin{remark}
Every term in \eqref{eq:fullmodel} is $C^\infty$ on the positive orthant
$\{x \in \mathbb{R}^5 : x_i > 0\}$, provided $\theta_i > 0$. This ensures that
standard tools for equilibrium existence (Brouwer/Schauder fixed point), local
stability (Jacobian linearization), and bifurcation theory apply directly without
the complications that would arise from the piecewise formulation of the adaptive
production term.
\end{remark}

\begin{remark}[Reduction to homeostasis model]
Setting $w_i = 0$ for all $i$ removes all Wu rebellion terms and yields a
homeostasis-focused model. The analysis of equilibria and their stability under
this reduction provides a baseline against which the pathological dynamics
introduced by Wu terms can be compared.
\end{remark}

\end{document}